% GENERATED FILE. Edit canonical lessons, then run npm run generate:books.
% canonical-source-hash: fnv1a64:4e7f5aaf7be43ae3
% canonical-lessons: SA-C429-udghatayati, SA-C429-ghatika, SA-C429-avasathah, SA-C429-argalah, SA-C429-avakrayah

\chapter{Opening Hours and Rent}
\label{ch:sa-opening-hours-and-rent}
\hlchaptermodality{429}

\begin{chapteropening}
\textbf{By the end of this chapter:} \emph{I can ask when a place opens and what a room costs.}

Complete the last lesson of chapter 429: I can ask when a place opens and what a room costs.
\end{chapteropening}

\section[\texorpdfstring{\sk{उद्घाटयति} (udghāṭayati)}{उद्घाटयति (udghāṭayati)}]{\sk{उद्घाटयति} (udghāṭayati) — opens; unlocks}

\label{lesson:SA-C429-udghatayati}



\begin{quote}
Before the new one: say the Sanskrit for goes out to meet, then the Sanskrit for sees a guest off.
\end{quote}

\subsection*{You'll want to know: \sk{उद्घाटयति}}
\textbf{\sk{उद्घाटयति}} — \emph{udghāṭayati} — \textquotedblleft{}opens; unlocks\textquotedblright{}.

\begin{quote}
\sk{आपणिकः} \sk{प्रातः} \sk{आपणम्} \sk{उद्घाटयति।}
\end{quote}

\begin{quote}
\emph{āpaṇikaḥ prātaḥ āpaṇam udghāṭayati.}
\end{quote}

\begin{quote}
The shopkeeper opens the shop in the morning.
\end{quote}

Past (he …): \sk{सः} \sk{आपणम्} \sk{उद्घाटितवान्}, \emph{saḥ āpaṇam udghāṭitavān}.

\begin{cousinweb}
\textbf{\sk{उद्घाटयति}} (\emph{udghāṭayati}) is built from \emph{ud}, \textquotedblleft{}up, out\textquotedblright{}, and the causative of the root √ghaṭ, which also gives \emph{ghaṭate}, \textquotedblleft{}happens\textquotedblright{}; MW gives \emph{ud-ghāṭayati} as to open, unlock, unfasten, with the participle \emph{udghāṭita} (MW).
\end{cousinweb}

\begin{grammarlens}[title={What you've built}]
One more word for getting around.
\end{grammarlens}

\subsection*{Your turn}
\begin{itemize}
\raggedright
  \item \emph{Say it:} \emph{udghāṭayati}
  \item \emph{Say it:} \emph{udghāṭayati}, once more
  \item \emph{Say it:} \emph{āpaṇikaḥ prātaḥ āpaṇam udghāṭayati.} — The shopkeeper opens the shop in the morning.
  \item \emph{Recall:} say the Sanskrit for goes out to meet, then the Sanskrit for sees a guest off, then say \emph{udghāṭayati} again
\end{itemize}

\begin{tcolorbox}[breakable,colback=teal!4,colframe=teal!35!black,title={Before you move on}]
What does \sk{उद्घाटयति} mean? (\textquotedblleft{}Opens; unlocks\textquotedblright{}.) Say it once more.
\end{tcolorbox}
\section[\texorpdfstring{\sk{घटिका} (ghaṭikā)}{घटिका (ghaṭikā)}]{\sk{घटिका} (ghaṭikā) — a clock; a ghatika of twenty-four minutes}

\label{lesson:SA-C429-ghatika}



\begin{quote}
Before the new one: say the Sanskrit for sees a guest off, then the Sanskrit for unlocks.
\end{quote}

\subsection*{You'll want to know: \sk{घटिका}}
\textbf{\sk{घटिका}} — \emph{ghaṭikā} — \textquotedblleft{}a clock; a ghatika of twenty-four minutes\textquotedblright{}.

\begin{quote}
\sk{आपणस्य} \sk{भित्तौ} \sk{घटिका} \sk{अस्ति।}
\end{quote}

\begin{quote}
\emph{āpaṇasya bhittau ghaṭikā asti.}
\end{quote}

\begin{quote}
There is a clock on the wall of the shop.
\end{quote}

A feminine noun.

\begin{cousinweb}
\textbf{\sk{घटिका}} (\emph{ghaṭikā}) is built on \emph{ghaṭī}, \textquotedblleft{}a small water-jar\textquotedblright{}: MW gives \emph{ghaṭikā} as a water-jar, as a period of twenty-four minutes once measured by such a jar, and as the Indian gong clock; a clock or a watch on the wall is that sense in modern use (MW).
\end{cousinweb}

\begin{grammarlens}[title={What you've built}]
One more word for getting around.
\end{grammarlens}

\subsection*{Your turn}
\begin{itemize}
\raggedright
  \item \emph{Say it:} \emph{ghaṭikā}
  \item \emph{Say it:} \emph{ghaṭikā}, once more
  \item \emph{Say it:} \emph{āpaṇasya bhittau ghaṭikā asti.} — There is a clock on the wall of the shop.
  \item \emph{Recall:} say the Sanskrit for sees a guest off, then the Sanskrit for unlocks, then say \emph{ghaṭikā} again
\end{itemize}

\begin{tcolorbox}[breakable,colback=teal!4,colframe=teal!35!black,title={Before you move on}]
What does \sk{घटिका} mean? (\textquotedblleft{}A clock; a ghatika of twenty-four minutes\textquotedblright{}.) Say it once more.
\end{tcolorbox}
\section[\texorpdfstring{\sk{आवसथः} (āvasathaḥ)}{आवसथः (āvasathaḥ)}]{\sk{आवसथः} (āvasathaḥ) — a travellers' lodge}

\label{lesson:SA-C429-avasathah}



\begin{quote}
Before the new one: say the Sanskrit for unlocks, then the Sanskrit for a clock.
\end{quote}

\subsection*{You'll want to know: \sk{आवसथः}}
\textbf{\sk{आवसथः}} — \emph{āvasathaḥ} — \textquotedblleft{}a travellers' lodge\textquotedblright{}.

\begin{quote}
\sk{आवसथे} \sk{पथिकाः} \sk{रात्रिं} \sk{वसन्ति।}
\end{quote}

\begin{quote}
\emph{āvasathe pathikāḥ rātriṃ vasanti.}
\end{quote}

\begin{quote}
Wayfarers spend the night at the lodge.
\end{quote}

A masculine noun.

\begin{cousinweb}
\textbf{\sk{आवसथः}} (\emph{āvasathaḥ}) is built from \emph{ā} and the root √vas, \textquotedblleft{}to dwell\textquotedblright{}, which also gives \emph{vasati}, \textquotedblleft{}dwells\textquotedblright{}; MW gives \emph{āvasatha} as a dwelling-place, an abode, a night's lodging (MW).
\end{cousinweb}

\begin{grammarlens}[title={What you've built}]
One more word for getting around.
\end{grammarlens}

\subsection*{Your turn}
\begin{itemize}
\raggedright
  \item \emph{Say it:} \emph{āvasathaḥ}
  \item \emph{Say it:} \emph{āvasathaḥ}, once more
  \item \emph{Say it:} \emph{āvasathe pathikāḥ rātriṃ vasanti.} — Wayfarers spend the night at the lodge.
  \item \emph{Recall:} say the Sanskrit for unlocks, then the Sanskrit for a clock, then say \emph{āvasathaḥ} again
\end{itemize}

\begin{tcolorbox}[breakable,colback=teal!4,colframe=teal!35!black,title={Before you move on}]
What does \sk{आवसथः} mean? (\textquotedblleft{}A travellers' lodge\textquotedblright{}.) Say it once more.
\end{tcolorbox}
\section[\texorpdfstring{\sk{अर्गलः} (argalaḥ)}{अर्गलः (argalaḥ)}]{\sk{अर्गलः} (argalaḥ) — a door bolt}

\label{lesson:SA-C429-argalah}



\begin{quote}
Before the new one: say the Sanskrit for a clock, then the Sanskrit for a travellers' lodge.
\end{quote}

\subsection*{You'll want to know: \sk{अर्गलः}}
\textbf{\sk{अर्गलः}} — \emph{argalaḥ} — \textquotedblleft{}a door bolt\textquotedblright{}.

\begin{quote}
\sk{रात्रौ} \sk{अहं} \sk{द्वारे} \sk{अर्गलं} \sk{स्थापयामि।}
\end{quote}

\begin{quote}
\emph{rātrau ahaṃ dvāre argalaṃ sthāpayāmi.}
\end{quote}

\begin{quote}
At night I put the bolt on the door.
\end{quote}

A masculine noun.

\begin{cousinweb}
\textbf{\sk{अर्गलः}} (\emph{argalaḥ}) is the word MW gives for a wooden bolt or pin that fastens a door, and for a bar or check; MW gives no derivation, so its origin is uncertain (MW).
\end{cousinweb}

\begin{grammarlens}[title={What you've built}]
One more word for getting around.
\end{grammarlens}

\subsection*{Your turn}
\begin{itemize}
\raggedright
  \item \emph{Say it:} \emph{argalaḥ}
  \item \emph{Say it:} \emph{argalaḥ}, once more
  \item \emph{Say it:} \emph{rātrau ahaṃ dvāre argalaṃ sthāpayāmi.} — At night I put the bolt on the door.
  \item \emph{Recall:} say the Sanskrit for a clock, then the Sanskrit for a travellers' lodge, then say \emph{argalaḥ} again
\end{itemize}

\begin{tcolorbox}[breakable,colback=teal!4,colframe=teal!35!black,title={Before you move on}]
What does \sk{अर्गलः} mean? (\textquotedblleft{}A door bolt\textquotedblright{}.) Say it once more.
\end{tcolorbox}
\section[\texorpdfstring{\sk{अवक्रयः} (avakrayaḥ)}{अवक्रयः (avakrayaḥ)}]{\sk{अवक्रयः} (avakrayaḥ) — rent}

\label{lesson:SA-C429-avakrayah}



\begin{quote}
Before the new one: say the Sanskrit for a travellers' lodge, then the Sanskrit for a door bolt.
\end{quote}

\subsection*{You'll want to know: \sk{अवक्रयः}}
\textbf{\sk{अवक्रयः}} — \emph{avakrayaḥ} — \textquotedblleft{}rent\textquotedblright{}.

\begin{quote}
\sk{प्रकोष्ठस्य} \sk{अवक्रयः} \sk{कियत्} \sk{अस्ति}?
\end{quote}

\begin{quote}
\emph{prakoṣṭhasya avakrayaḥ kiyat asti?}
\end{quote}

\begin{quote}
How much is the rent for the room?
\end{quote}

A masculine noun.

\begin{cousinweb}
\textbf{\sk{अवक्रयः}} (\emph{avakrayaḥ}) is built from \emph{ava} and the root √krī, \textquotedblleft{}to buy\textquotedblright{}, which also gives \emph{krīṇāti}, \textquotedblleft{}buys\textquotedblright{}; MW gives \emph{avakraya} as letting out to hire and as rent (MW).
\end{cousinweb}

\begin{grammarlens}[title={What you've built}]
One more word for getting around.
\end{grammarlens}

\subsection*{Your turn}
\begin{itemize}
\raggedright
  \item \emph{Say it:} \emph{avakrayaḥ}
  \item \emph{Say it:} \emph{avakrayaḥ}, once more
  \item \emph{Say it:} \emph{prakoṣṭhasya avakrayaḥ kiyat asti?} — How much is the rent for the room?
  \item \emph{Recall:} say the Sanskrit for a travellers' lodge, then the Sanskrit for a door bolt, then say \emph{avakrayaḥ} again
\end{itemize}

\begin{tcolorbox}[breakable,colback=teal!4,colframe=teal!35!black,title={Before you move on}]
What does \sk{अवक्रयः} mean? (\textquotedblleft{}Rent\textquotedblright{}.) Say it once more.
\end{tcolorbox}
